\documentclass[11pt]{article}
\usepackage{docstyle}
\renewcommand\sheetname{Handout 3 \textperiodcentered{} Quadratics and the discriminant}

\begin{document}
\sheethead{Lesson 3 \textperiodcentered{} Quadratics and the Discriminant}
  {NCEA Level 2 Mathematics \textperiodcentered{} Algebra}{Handout}

\begin{keybox}[title={The one idea: get = 0 first}]
A product is zero only when one of its factors is zero. So $(x-3)(x+2) = 0$ gives $x = 3$ or $x = -2$, but $(x-3)(x+2) = 6$ tells you nothing about either bracket.
\textbf{Rearrange to $ax^{2} + bx + c = 0$ before you do anything else.} Then three tools work: factorise, the formula, and the \Tdisc{} $b^{2} - 4ac$, which counts the solutions before you find them.
\end{keybox}

% =================================================================== 1
\section{Solving a quadratic}\label{h3:solve}
\begin{keybox}[title={Write it in this order}]
\begin{enumerate}
\item Rearrange to $ax^{2} + bx + c = 0$. Divide through by a common factor if there is one.
\item Try to factorise. If it factorises, set each bracket to zero.
\item If it does not, use the formula (on the formula sheet), with $a$, $b$, $c$ \textbf{in brackets}:
\[ x = \frac{-b \pm \sqrt{b^{2} - 4ac}}{2a} \]
\item Give exact answers (surds) when asked; otherwise round at the end.
\end{enumerate}
\end{keybox}

\begin{multicols}{2}
\textbf{Example 1.} Solve $x(x+2) = 15$.
\begin{align*}
x^{2} + 2x - 15 &= 0\\
(x+5)(x-3) &= 0\\
x = \mathbf{-5} \text{ or } x &= \mathbf{3}
\end{align*}
\mtag{u}
Not $x = 15$: the right side was not zero.

\textbf{Example 2.} Solve $2x^{2} - 5x - 1 = 0$.\par
$a = 2$, $b = -5$, $c = -1$:
\[ x = \frac{-(-5) \pm \sqrt{(-5)^{2} - 4(2)(-1)}}{2(2)} = \frac{5 \pm \sqrt{33}}{4} \]
$x = \mathbf{2.686}$ or $x = \mathbf{-0.186}$ \mtag{u}
\columnbreak

\textbf{Example 3.} Solve $x^{2} - 6x + 2 = 0$ by completing the square.
\begin{align*}
x^{2} - 6x &= -2\\
(x - 3)^{2} - 9 &= -2\\
(x - 3)^{2} &= 7\\
x &= \mathbf{3 \pm \sqrt{7}}
\end{align*}
\mtag{u}
Halve the $x$ coefficient ($-6 \to -3$) and take away its square.

\textbf{Example 4.} Solve $x^{4} - 13x^{2} + 36 = 0$.\par
Let $u = x^{2}$: $u^{2} - 13u + 36 = (u-4)(u-9) = 0$, so $x^{2} = 4$ or $x^{2} = 9$: $x = \mathbf{\pm 2}$ or $\mathbf{\pm 3}$. \mtag{r}
\end{multicols}

% =================================================================== 2
\section{From the roots back to the equation}\label{h3:roots}
A root $x = r$ comes from the factor $(x - r)$; a root $x = \dfrac{p}{q}$ comes from the factor $(qx - p)$, which keeps the coefficients whole.
\begin{multicols}{2}
\textbf{Example 5.} A quadratic has roots $\dfrac{3}{5}$ and $-2$. Find it with whole-number coefficients.
\begin{align*}
x = \tfrac{3}{5} &\Rightarrow (5x - 3)\\
x = -2 &\Rightarrow (x + 2)\\
(5x-3)(x+2) &= \mathbf{5x^{2} + 7x - 6 = 0}
\end{align*}
\mtag{r}
\columnbreak

\textbf{Example 6.} A parabola crosses the $x$-axis at $-1$ and $5$ and passes through $(0, -10)$.
\begin{align*}
y &= a(x+1)(x-5)\\
-10 &= a(1)(-5) \;\Rightarrow\; a = 2\\
y &= \mathbf{2x^{2} - 8x - 10}
\end{align*}
\mtag{r}
\end{multicols}

% =================================================================== 3
\section{The discriminant}\label{h3:disc}
\[ \Delta = b^{2} - 4ac \]
\begin{center}
\setlength\figunit{4.5mm}
\begin{tabular}{@{}c c c@{}}
\DiscPanel{-1} & \DiscPanel{0} & \DiscPanel{1}\\[1mm]
$b^{2} - 4ac > 0$ & $b^{2} - 4ac = 0$ & $b^{2} - 4ac < 0$\\
two different real roots & one repeated root (equal roots) & no real roots
\end{tabular}
\end{center}
``Real roots'' (without ``different'') means $b^{2} - 4ac \geq 0$.

\begin{keybox}[title={Write it in this order (find $k$)}]
\begin{enumerate}
\item Read the condition and choose the sign: $> 0$, $= 0$, $< 0$ or $\geq 0$.
\item Write $a$, $b$, $c$, each \textbf{in brackets}, with $k$ wherever it appears.
\item Substitute into $b^{2} - 4ac$ and solve the equation or inequality. Dividing an inequality by a negative number \textbf{turns it round}.
\end{enumerate}
\end{keybox}

\begin{multicols}{2}
\textbf{Example 7.} $x^{2} + kx + 16 = 0$ has equal roots. Find $k$.
\[ k^{2} - 4(1)(16) = 0 \;\Rightarrow\; k = \mathbf{\pm 8} \]
\mtag{u}

\textbf{Example 8.} $2x^{2} + 6x + p = 0$ has two different real roots. Find $p$.
\[ 36 - 8p > 0 \;\Rightarrow\; -8p > -36 \;\Rightarrow\; \mathbf{p < 4.5} \]
\mtag{r}
\columnbreak

\textbf{Example 9.} $kx^{2} + 4x + (k - 3) = 0$ has equal roots. Find $k$.
\begin{align*}
4^{2} - 4(k)(k-3) &= 0\\
k^{2} - 3k - 4 &= 0\\
k = \mathbf{4} \text{ or } k &= \mathbf{-1}
\end{align*}
\mtag{r}
\end{multicols}

% =================================================================== 4
\section{Proving something about the roots}\label{h3:prove}
\textbf{Example 10.} Show that one solution of $x^{2} - 5kx + 4k^{2} = 0$ ($k \neq 0$) is four times the other.
\[ (x - k)(x - 4k) = 0 \;\Rightarrow\; x = k \text{ or } x = 4k, \text{ and } 4k = 4 \times k. \]
\mtag{r}
Or by the formula: $\Delta = 25k^{2} - 16k^{2} = 9k^{2}$, so $x = \dfrac{5k \pm 3k}{2} = 4k$ or $k$. \textbf{Finish with the sentence the question asks for.}

\textbf{Always real roots.} To show an equation has real roots for \textbf{every} $k$, write $\Delta$ as a square: $(k+2)^{2} - 8k = (k-2)^{2} \geq 0$.

% =================================================================== 5
\section{Versions the marker won't accept}\label{h3:wontaccept}
\begin{tabularx}{\textwidth}{@{}>{\raggedright}p{0.3\textwidth} >{\raggedright}p{0.26\textwidth} X@{}}
\toprule
\textbf{Question} & \textbf{Answer given} & \textbf{Why it falls short} \\ \midrule
Roots $-\frac{2}{3}$ and $4$; give $ax^{2} + bx + c = 0$ with whole numbers & $x^{2} - \frac{10}{3}x - \frac{8}{3} = 0$ & Correct but not whole numbers: Achievement evidence only. Multiply by 3: $3x^{2} - 10x - 8 = 0$. \tabularnewline
$y = 0$ has exactly one solution & $144 - 8k > 0$ & Wrong condition: one solution means $\Delta = 0$. \tabularnewline
Show one solution is twice the other & $x = k$ or $x = 2k$ & Add the conclusion: ``so one solution is twice the other''. \tabularnewline
\bottomrule
\end{tabularx}

% =================================================================== 6
\section{Ways to lose marks}\label{h3:lose}
From the wording of the marking schedules and from teaching observation; numbered for marking.
\begin{description}[leftmargin=10mm, labelwidth=9mm, labelsep=1mm, font=\normalfont]
\item[\textbf{\Lref{rootsign}}] \textbf{\Ltext{rootsign}.} A root of $4$ gives $(x - 4)$, a root of $\frac{2}{3}$ gives $(3x - 2)$.
\item[\textbf{\Lref{whole}}] \textbf{\Ltext{whole}.} Clear fractions before expanding, or multiply through at the end.
\item[\textbf{\Lref{notzero}}] \textbf{\Ltext{notzero}.} $x(x+3) = 10$ does not give $x = 10$.
\item[\textbf{\Lref{disccond}}] \textbf{\Ltext{disccond}.} Two different: $>0$. Equal or exactly one: $=0$. None: $<0$. Real: $\geq 0$.
\item[\textbf{\Lref{discbracket}}] \textbf{\Ltext{discbracket}.} $(-5)^{2} = 25$, but $-5^{2} = -25$. Brackets round $b$, and round $k$ terms.
\end{description}

% =================================================================== 7
\section{Quick review}\label{h3:review}
\renewcommand\arraystretch{1.45}
\begin{minipage}[t]{0.47\textwidth}\vspace{0pt}
\textbf{Factorising (Lesson 1)}\par\smallskip
\begin{tabular}{@{}l l@{}}
\toprule
$x^{2} + 5x + 6$ & $(x+2)(x+3)$ \\
$x^{2} - 9$ & $(x-3)(x+3)$ \\
$2x^{2} + 5x + 2$ & $(2x+1)(x+2)$ \\
$3x^{2} - 6x$ & $3x(x-2)$ \\
\bottomrule
\end{tabular}
\end{minipage}\hfill
\begin{minipage}[t]{0.47\textwidth}\vspace{0pt}
\textbf{Linear inequalities}\par\smallskip
\begin{tabular}{@{}l l@{}}
\toprule
$3x - 2 > 7$ & $x > 3$ \\
$12 - 4k < 0$ & $-4k < -12$, so $k > 3$ \\
$-2p \geq 8$ & $p \leq -4$ \\
\bottomrule
\end{tabular}\par\smallskip
Multiply or divide by a negative: turn the sign round.
\end{minipage}
\end{document}
